% ---- pg-510 （书页 498） ----

\begin{mvframe}[显示板 11.2(续)]
{\small
{\centering Probability modeled is country $=2$.\\[2pt] Model Fit Statistics\par}
\vspace{6pt}
{\centering
\begin{tabular}{lcc}
 & Intercept & Intercept\\
 & Only & and\\
Criterion & & Covariates\\
AIC & 140.629 & 46.674\\
SC & 143.235 & 57.094\\
-2 Log L & 138.629 & 38.674\\
\end{tabular}\par}
\vspace{10pt}
{\centering Testing Global Null Hypothesis: BETA $=0$\par}
\vspace{4pt}
{\centering
\begin{tabular}{lccc}
Test & Chi-Square & DF & Pr $>$ ChiSq\\
\hline
Likelihood Ratio & 99.9557 & 3 & $<.0001$\\
Wald & 19.4435 & 3 & 0.0002\\
\hline
\end{tabular}\par}
\vspace{10pt}
{\centering The LOGISTIC Procedure\par}
\vspace{6pt}
{\centering Analysis of Maximum Likelihood Estimates\par}
\vspace{4pt}
{\centering
\begin{tabular}{lcccccc}
Parameter & DF & Estimate & \makecell{Standard\\Error} & \makecell{Wald\\Chi-Square} & Pr $>$ ChiSq & Exp (Est)\\
\hline
Intercept & 1 & 3.5054 & 6.3935 & 0.3006 & 0.5835 & 33.293\\
gender & 1 & 0.2816 & 0.8338 & 0.1140 & 0.7356 & 1.325\\
freshwater & 1 & 0.1264 & 0.0357 & 12.5423 & 0.0004 & 1.135\\
marine & 1 & $-0.0486$ & 0.0146 & 11.1460 & 0.0008 & 0.953\\
\hline
\end{tabular}\par}
}
\end{mvframe}

我们还可以在例 11.17 中去掉性别, 对淡水生长和海水生长两个预测变量作逻辑斯蒂回归分析. 但是, 与具有相同协方差阵的正态分布相比, 逻辑斯蒂分类的效率要低得多(见文献[17]).

\subsection{二项分布响应的逻辑斯蒂回归}

现在我们考虑稍微一般一些的情况. 假定有些观测是在相同的协变量 $\mathbf{z}_j$ 下得到的, 并假定总共有 $m$ 个数据集, 其中每个数据集内的预测变量相同. 当在预测变量 $\mathbf{z}_j$ 下作 $n_j$ 次独立试验, 响应 $Y_j$(译者注: $n_j$ 次试验中响应为 1 的频数)有二项分布, 成功概率为 $p(\mathbf{z}_j)=P(\text{成功}\mid \mathbf{z}_j)$.

由于 $Y_j$ 是独立的, $j=1,\cdots,m$, 因此似然函数表为乘积:
\phantomsection\label{eq:11-78}
\begin{equation}
L(\beta_0,\beta_1,\cdots,\beta_r)=\prod_{j=1}^{m}\binom{n_j}{y_j}p_j^{y_j}(\mathbf{z}_j)[1-p(\mathbf{z}_j)]^{n_j-y_j}
\tag{11-78}
\end{equation}
其中概率 $p(\mathbf{z}_j)$ 遵从逻辑模型(11-72).

极大似然估计 $\hat{\boldsymbol{\beta}}$ 没有封闭的函数形式, 只能得到数值解. 当样本的总容量足够大时, 渐近的协方差矩阵为
\phantomsection\label{eq:11-79}
\begin{equation}
\widehat{\mathrm{Cov}}(\hat{\boldsymbol{\beta}})\approx\left[\sum_{j=1}^{m}n_j\hat{p}(\mathbf{z}_j)(1-\hat{p}(\mathbf{z}_j))\mathbf{z}_j\mathbf{z}_j'\right]^{-1}
\tag{11-79}
\end{equation}
其中第 $i$ 个对角元为 $\hat{\beta}_{j+1}$ 的方差估计. 其平方根为标准误差 $\mathrm{SE}(\hat{\beta}_{j+1})$ 的大样本估计.

可以证明, 概率估计 $\hat{p}(\mathbf{z}_j)$ 的方差的大样本估计为
